\scanpage{60}
\mspage{9}
% archive sheet number: 60

% slanted
\margin{on utilise ici $W(2\,\mathrm{bis})$, pas $W(2)$}

prouver que les fibres de $p$ sont $W$-asphériques. Or les fibres sont les $Z_\tau=\underline{\mathrm{Ch}}_\tau(S;s,-)$, qui sont des catégories \uline{contractiles} (i.e. $Z_\tau\unsure{\text{---}|\longrightarrow}e$ est un homotopisme), donc $W$-asphériques par \textcircled{3} (cf \unsure{cor}. 3 page 3).

D'autre part, \deleted{voyons que} si on admet que $\underline{\mathrm{Ch}}(S)\xrightarrow{p_S}S\times S$ est $W$-lisse (et $W$-propre)\deleted{\unsure{et}}, donc $\underline{\mathrm{Ch}}(S)\xrightarrow{\sigma_S}S$ aussi, étant à fibres $W$-asphériques comme on vient de voir, il s'ensuit que $\sigma_S\in W$, d'où $S\xrightarrow{i_S}\underline{\mathrm{Ch}}(S)$ aussi dans $W$ (puisque $\deleted{\unsure{p}}\sigma_S\,i_S=\mathrm{id}_S$).

% slanted
\margin{\unsure{id}}

Donc il reste à prouver seulement que
\[
\underline{\mathrm{Ch}}(S)\xrightarrow{p_S}S\times S \tag*{(a)}
\]
est $W$-\deleted{\unsure{lisse}}propre et $W$-lisse. Il en résultera que
\[
\underline{\mathrm{Ch}}(S;s_0,-)\xrightarrow{\beta}S
\]
est lui aussi $W$-propre et $W$-lisse, car déduit de $p_S$ par le chgt de base
\[
S\longrightarrow S\times S\qquad s\longmapsto(s_0,s).
\]

Il faut ensuite voir que \deleted{\unsure{$f$}} $W$-lisse et $W$-propre implique la propriété de $W$-fibration \uline{sous sa forme la plus forte} a) (alors que je prévois qu'a priori on ne l'aura que sous les formes faibles $a_1,a_2$ -- pas même sous la forme $a'$). C'est ici donc que l'axiome \textcircled{4} des $W$-fibrations sera \unsure{appréciable|apprécié} !

% slanted
\margin{\unsure{sous} la forme affaiblie $\unsure{W(4a'')|4a'')}$ (p. 13)}

Voici l'argument pour la lissité de \unsure{(a)|cas}, dans loc. cit p. 161 ff. Il faut prouver que \deleted{la catégorie}\inserted{pour} \deleted{\unsure{$c\in\mathrm{Ob}\,\underline{\mathrm{Ch}}$}} $c\in\mathrm{Ob}\,\underline{\mathrm{Ch}}(S)$, \deleted{\unsure{$(a,b)\in\mathrm{Ob}$}} et une flèche $u\in\mathrm{Fl}(S\times S)$, une certaine catégorie du ``relèvement de $u$ de source $c$'', $C=C(c,u,\unsure{v|u})$, est $W$-asphérique. Pour ceci, on la compare à une catégorie \inserted{$C'$} à objet final, donc \inserted{contractile et à fortiori} $W$-asphérique, à l'aide de \deleted{\unsure{du}} foncteurs

% slanted
\margin{En fait, on prouve que $C(c,u,v)$ est \uline{contractile}, donc elle est $W$-asphérique}

\diagram{\begin{tikzcd}[ampersand replacement=\&,column sep=large]
C\arrow[r,shift left=.65ex,"\varphi",dash] \& C'\arrow[l,shift left=.65ex,"\psi"]
\end{tikzcd}}
, et des homotopies $\varphi\psi\sim\mathrm{id}_{C'}$ (donc, puisque $C'$ est contractile)
